% !TeX root = main.tex
\chapter{Metodologia de Pesquisa}
\label{cap:metodologia}

Este capítulo descreve os procedimentos metodológicos adotados para a concepção, o desenvolvimento e a avaliação experimental do artefato proposto. O trabalho apoia-se nas diretrizes da \textit{Design Science Research} (DSR), método voltado à resolução de problemas organizacionais mediante a criação e avaliação rigorosa de artefatos computacionais \cite{hevner2004,peffers2007}.

% ============================================================
\section{Caracterização Metodológica}
\label{sec:dsr-caracterizacao}
% ============================================================

\subsection{Design Science Research como Paradigma}
\label{subsec:dsr-paradigma}

A pesquisa em computação aplicada comporta abordagens explicativas e prescritivas. Enquanto as ciências comportamentais focam na observação de fenômenos, a DSR orienta-se pela produção de conhecimento por meio da construção e da validação empírica de artefatos --- constructos, modelos, métodos e instanciações \cite{hevner2004}. 

A gestão de visitas em imóveis desocupados apresenta gargalos operacionais e vulnerabilidades que decorrem de métodos manuais (chaves físicas e senhas estáticas) e da ausência de internet local nos imóveis. A resposta a essa demanda requer a especificação de modelos formais e a implementação de protótipos funcionais. A condução do estudo segue o modelo de processo de seis etapas estabelecido por \citeonline{peffers2007}, estruturado conforme o Quadro~\ref{quadro:mapeamento-dsr}.

\begin{table}[htbp]
\centering
\footnotesize
\caption{Mapeamento das etapas da Design Science Research no projeto.}
\label{quadro:mapeamento-dsr}
\begin{tabularx}{\textwidth}{c p{3.8cm} X p{3.0cm}}
\toprule
\textbf{Fase} & \textbf{Etapa \cite{peffers2007}} & \textbf{Atividade Realizada} & \textbf{Localização} \\
\midrule
1 & Identificação do Problema & Diagnóstico das falhas de segurança e logística em imóveis vagos sem conectividade. & Seção~\ref{sec:problema-contexto} e Cap.~\ref{cap:introducao} \\
\addlinespace
2 & Definição dos Objetivos & Definição dos requisitos de projeto, pilares arquiteturais e metas quantitativas. & Seção~\ref{subsec:pilares-artefato} e Seção~\ref{sec:metricas-avaliacao} \\
\addlinespace
3 & Projeto e Desenvolvimento & Concepção do modelo HABAC$\alpha$, esquema Ed25519, firmware ESP32 e API FastAPI. & Cap.~\ref{cap:resultados} \\
\addlinespace
4 & Demonstração & Execução do fluxo nominal, testes negativos de ataque e revogação em bancada virtual. & Seção~\ref{sec:cenarios-teste} e Cap.~\ref{cap:resultados} \\
\addlinespace
5 & Avaliação Experimental & Coleta de latência, memória e eficácia de bloqueio em 50 repetições com análise estatística. & Seção~\ref{sec:metricas-avaliacao} e Cap.~\ref{cap:resultados} \\
\addlinespace
6 & Comunicação & Redação acadêmica, sistematização de diagramas e discussão comparativa dos dados. & Monografia \\
\bottomrule
\end{tabularx}
\fonte{Elaborado pelo autor (2026).}
\end{table}

% ============================================================
\section{Contexto e Formalização do Problema}
\label{sec:problema-contexto}
% ============================================================

\subsection{Cenário de Gestão de Visitas e Análise dos Métodos Vigentes}
\label{subsec:cenario-imobiliario}

Imobiliárias administram carteiras distribuídas de unidades desocupadas cuja comercialização depende de corretores parceiros e prestadores de serviços. A operação caracteriza-se por três restrições: (i) ausência de conexão de rede local (Wi-Fi ou cabeada), visto que custos de banda larga não são mantidos na vacância; (ii) inexistência de fiscalização física no local; e (iii) elevada rotatividade de profissionais que necessitam de janelas de entrada temporárias e pontuais.

Conforme analisado no Quadro~\ref{quadro:comparativo-metodos} (Seção~\ref{subsec:limitacoes-tradicionais}), os métodos convencionais mostram-se inadequados: chaves metálicas causam atrito logístico e risco de cópias; cofres mecânicos empregam senhas estáticas compartilhadas sem rastreabilidade; e fechaduras comerciais conectadas dependem de Wi-Fi local ininterrupto, tornando-se inoperantes em imóveis desocupados.

\subsection{Questão de Pesquisa e Pilares do Artefato}
\label{subsec:questao-pesquisa}

A incapacidade dos mecanismos atuais de conciliar operação desconectada, controle contextual e auditoria orienta a questão de pesquisa:

\begin{quote}
\textit{Como projetar e avaliar uma arquitetura de controle de acesso inteligente baseada em IoT capaz de gerenciar, de forma segura, rastreável e orientada a atributos (ABAC), o acesso temporário de múltiplos corretores parceiros a imóveis desocupados desprovidos de conexão à internet?}
\end{quote}

\label{subsec:pilares-artefato}
Para orientar a construção do artefato, definiram-se quatro pilares técnicos:
\begin{enumerate}
    \item \textbf{Independência de rede local (\textit{Zero Local Infrastructure}):} Validação matemática de credenciais na borda sem consulta a servidores remotos e sem modems dedicados na fechadura.
    \item \textbf{Autenticidade e não-repúdio:} Delegação assinada por autoridade central e vinculada ao dispositivo móvel habilitado do corretor.
    \item \textbf{Autorização orientada a atributos (ABAC):} Liberação condicionada a atributos de sujeito, objeto, ação e ambiente avaliados em tempo de execução.
    \item \textbf{Auditoria assíncrona:} Registro inviolável de acessos e sincronização com a nuvem utilizando o smartphone do corretor como portador de dados (\textit{data mule}).
\end{enumerate}

% ============================================================
\section{Procedimentos e Instrumentos de Avaliação}
\label{sec:procedimentos-avaliacao}
% ============================================================

\subsection{Ambiente de Simulação Virtual}
\label{subsec:justificativa-wokwi}

A avaliação do artefato adota a simulação virtual no Wokwi Simulator \cite{simperf2024}. A opção justifica-se por três fatores: (i) reprodutibilidade científica integral dos ensaios de temporização em ambiente controlado; (ii) isolamento de variáveis espúrias mecânicas ou térmicas de protoboards na medição da computação criptográfica; e (iii) viabilidade experimental de testes de estresse sem degradação física de atuadores solenoide.

\subsection{Planejamento dos Cenários de Demonstração}
\label{sec:cenarios-teste}

A validação funcional divide-se em três cenários planejados:
\label{subsec:cenario-1}
\textbf{Cenário 1 (Fluxo nominal):} Avalia a solicitação de token por corretor cadastrado para imóvel liberado em horário comercial. Critério de sucesso: emissão do token, validação no microcontrolador em menos de 10~ms, acionamento do relé por 5 segundos e entrega do recibo de telemetria via MQTT sob QoS 1.

\label{subsec:cenario-2}
\textbf{Cenário 2 (Ensaios negativos e ataques):} Avalia a capacidade de rejeição local do firmware frente a credenciais anômalas, de acordo com o planejamento do Quadro~\ref{quadro:cenario-2-testes}. Critério de aceitação: índice de 100\% de bloqueio efetivo sem concessões indevidas.

\begin{table}[htbp]
\centering
\footnotesize
\caption{Planejamento dos ensaios de segurança do Cenário 2.}
\label{quadro:cenario-2-testes}
\begin{tabularx}{\textwidth}{c p{3.2cm} X p{3.6cm}}
\toprule
\textbf{Caso} & \textbf{Condição Testada} & \textbf{Estímulo Aplicado} & \textbf{Comportamento Esperado} \\
\midrule
CT-2.1 & Token Expirado & Credencial autêntica com timestamp posterior ao campo \texttt{exp}. & Rejeição por TTL estourado e LED vermelho. \\
CT-2.2 & Imóvel Incorreto & Token emitido para a unidade A apresentado na unidade B. & Bloqueio por divergência no campo \texttt{dev}. \\
CT-2.3 & Assinatura Corrompida & Adulteração intencional de 1 bit no conteúdo do payload JSON. & Falha na verificação Ed25519 e descarte. \\
CT-2.4 & Ataque de Repetição & Reenvio de credencial já consumida em momento anterior. & Bloqueio por identificador \texttt{jti} no buffer NVS. \\
\bottomrule
\end{tabularx}
\fonte{Elaborado pelo autor (2026).}
\end{table}

\label{subsec:cenario-3}
\textbf{Cenário 3 (Revogação e risco residual):} Avalia a resposta do sistema diante da suspensão cadastral de um corretor. O plano verifica a imediata recusa de novas emissões pelo servidor (HTTP 403) e a contenção no nó desconectado por tempo de vida estrito (\textit{Time-To-Live} de 15 minutos) para tokens emitidos previamente.

\subsection{Métricas Quantitativas e Tratamento Estatístico}
\label{sec:metricas-avaliacao}
\label{subsec:variaveis-metricas}

Para a avaliação empírica do artefato, definiram-se oito métricas quantitativas com base na literatura de IoT \cite{ameer2022,vardakis2024,simperf2024}, sintetizadas no Quadro~\ref{quadro:metricas-avaliacao}.

\begin{table}[htbp]
\centering
\footnotesize
\caption{Métricas quantitativas planejadas para a avaliação experimental.}
\label{quadro:metricas-avaliacao}
\begin{tabularx}{\textwidth}{c p{4.2cm} c X p{2.2cm}}
\toprule
\textbf{\#} & \textbf{Métrica} & \textbf{Unid.} & \textbf{Instrumento e Coleta} & \textbf{Meta} \\
\midrule
M1 & Latência de Decisão ABAC & ms & Medição de tempo de processamento no PDP FastAPI. & $<$ 50 ms \\
M2 & Tempo de Assinatura Criptográfica & ms & Medição interna da biblioteca PyNaCl (Ed25519). & $<$ 5 ms \\
M3 & Verificação no ESP32 & ms & Temporização via função \texttt{micros()} no firmware. & $<$ 10 ms \\
M4 & Latência Ponta a Ponta & ms & Intervalo entre requisição móvel e acionamento do relé. & $<$ 500 ms \\
M5 & Consumo de Memória RAM Heap & KB & Amostragem via rotina nativa \texttt{ESP.getFreeHeap()}. & $<$ 120 KB \\
M6 & Ocupação de Memória Flash & \% & Relatório de vinculação estática do PlatformIO. & $<$ 35\% \\
M7 & Vazão de Mensagens MQTT & msg/s & Teste de estresse de publicações com Locust. & $>$ 100 msg/s \\
M8 & Taxa de Bloqueio de Ataques & \% & Percentual de rejeição nos testes CT-2.1 a CT-2.4. & 100\% \\
\bottomrule
\end{tabularx}
\fonte{Elaborado pelo autor (2026).}
\end{table}

\label{subsec:procedimento-experimental}
O procedimento experimental prevê 50 repetições para cada medição temporal, reduzindo perturbações estatísticas. Os dados a serem coletados nas baterias de testes serão tratados estatisticamente com as bibliotecas Pandas e SciPy, apurando média ($\mu$), desvio padrão ($\sigma$), valores extremos e os percentis $p50$, $p95$ e $p99$. Os resultados esperados e o detalhamento analítico do artefato constam no Capítulo~\ref{cap:resultados}.
